\documentclass{article}

\usepackage{neurips_2025}

\usepackage[utf8]{inputenc} % allow utf-8 input
\usepackage[T1]{fontenc}    % use 8-bit T1 fonts
\usepackage{hyperref}       % hyperlinks
\usepackage{url}            % simple URL typesetting
\usepackage{graphicx}       % resizebox
\usepackage{booktabs}       % professional-quality tables
\usepackage{amsfonts}       % blackboard math symbols
\usepackage{nicefrac}       % compact symbols for 1/2, etc.
\usepackage{microtype}      % microtypography
\usepackage{xcolor}         % colors
\usepackage{amsmath}        % math environments
\usepackage{algorithm}      % algorithm environment
\usepackage{algorithmic}    % algorithmic pseudocode

\title{Self-Evolving Game Agent with MCTS and Tool-Integration}

\author{%
  Jason Zhu\thanks{Equal contribution.} \\
  UC San Diego\\
  \texttt{zhuyuezx@ucsd.edu} \\
  \And
  Yinuo Zhou\footnotemark[1] \\
  UC San Diego\\
  \And
  Hongrui (Frank) Zhang\footnotemark[1] \\
  UC San Diego\\
  \And
  Yaran Yang\footnotemark[1] \\
  UC San Diego\\
  \And
  Zavier Zhao\footnotemark[1] \\
  UC San Diego\\
  \And
  Fengfei Yu\footnotemark[1] \\
  UC San Diego\\
}


\begin{document}


\maketitle


\begin{abstract}
% TODO: Write abstract based on the following key points from the presentation:
% - We develop a self-evolving LLM tool-creation agent for turn-based games
% - The agent interacts with game environments (Sokoban, TextWorld, Zork), collects traces, analyzes feedback, and generates pluggable heuristics
% - Uses MCTS as a foundational framework and also explores A* for games where MCTS struggles
% - LLM acts as an offline coding agent producing human-understandable code-based heuristics, in contrast to online LLM decision agents or deep neural network training
% - Generated heuristics show improvement even with low iterations and demonstrate transferability
% - Evaluated on Sokoban (A*), TextWorld benchmarks (MCTS), and Zork, showing significant improvements over baselines
\end{abstract}


\section{Introduction}
\label{sec:intro}

% Motivation and core problems:
% - Solving turn-based games as MDP problems has been an effective strategy (Go, etc.)
% - Monte Carlo Tree Search searches the game environment and finds promising paths algorithmically
% - AlphaZero-style MCTS trains deep neural networks for state value estimations, which is hard to set up and has massive entry cost
% - Earlier MCTS+LLM work focused on online, step-wise LLM guidance, which is costly and constrained by static local views
% - Emergence of LLM skills/MCP gave motivation for another direction: policy improvement with rule-based heuristics proposed autonomously by LLM

% Our contribution:
% - We developed a LLM tool-evolving agent for turn-based games that:
%   - Interacts with different game environments (Sokoban, TextWorld, etc.)
%   - Collects game traces, analyzes feedback, does task proposal, generates pluggable heuristics, improves upon successful validation, and iterates
%   - Mainly uses MCTS as a foundational framework, but also explores A* for games where MCTS struggles
%   - Generates heuristics for MCTS, with improvement seen even at low iterations
%   - LLM acts as an offline coding agent, while code-based heuristics run online efficiently
%   - Generated heuristics are human-understandable code artifacts instead of opaque weights

% Audience of interest:
% - LLM users and researchers: explores LLM tool creation and self-evolving agent paradigm
% - Task proposal: LLM's ability to learn from environment feedback and propose improvements autonomously
% - Game players: generating heuristics offline is less expensive and leads to efficient gameplay improvements


\section{Related Work}
\label{sec:related}

\subsection{A Survey of Monte Carlo Tree Search Methods}
% [2012] Cameron Browne et al.
% - Challenge: How to find optimal decisions in a given domain (complex game as MDP)?
% - Paper: Maps the problem to a search tree, performs 4 stages (selection, exploration, simulation, back-propagation) with tree policy to understand game states and identify optimal decisions
% - Exploration heuristics like UCT balance exploration and exploitation with simple data
% - Limitation: Lacks reasoning over environment data, constrained by limited information
% - Our learning: Solve complex games based on an MCTS framework
% - Our approach: Use LLM to enhance reasoning over environment and state information

\subsection{MC-DML: Monte-Carlo Planning with Dynamic Memory-Guided LLM}
% [ICLR 2025] Zijing Shi et al.
% - Challenge: Past MCTS+RL limitations, like uncertain exploration and lack of reasoning
% - Paper: Deploys LLM as a policy (action decision) maker in text games, enriches context with dynamic memory mechanism (in-track, cross-track memory)
% - Limitation: High cost, confusion with local memory management
% - Our learning: Utilize LLM's reasoning capabilities and memory management
% - Our approach: We generate heuristics used for many steps; we study summarized traces from a global point of view for tool generation

\subsection{Yunjue Agent: A Self-Evolving Agent System for Open-Ended Tasks}
% [2026] Haotian Li et al.
% - Challenge: LLM requires human supervision to solve hard problems
% - Paper: Multi-agent architecture that solves sequences of real-world problems by generating, synthesizing, optimizing, and reusing tools (python scripts) in parallel
% - Limitation: Focuses on deterministic problems, can benefit from targeted workflow
% - Our learning: Novel tool generation paradigm reduces supervision in our system
% - Our approach: Using LLM to analyze, generate, and improve MCTS or other heuristics as tools for validated agentic task proposals


\section{Method}
\label{sec:method}

\subsection{System Architecture}
% - RL framework not limited to MCTS
% - Modularize optimizable sections, allowing LLM creation and modification
% - Adaptable to all turn-based games
% - Codebase: https://github.com/zhuyuezx/CSE291A_Project

% Framework selection rationale:
% - A* is suitable for games with simple action space and easy-to-get heuristics but complex, irreversible state space (e.g., Sokoban)
% - MCTS is suitable for games with complex action space and hard-to-get heuristics (e.g., Go, Quoridor)

\subsection{Task Proposal Loop}
% - Input: Current tool code + environment feedback given current tools
% - Mutable source for improvement seed
% - Improvement ideas generated iteratively
% - Verified and tested before being incorporated
% - Task proposal quality is the most important fault point


\section{Experiments}
\label{sec:experiments}

\subsection{Sokoban (A* Optimization)}
% - Environment description: Sokoban puzzle game with boxes and goals on a grid
% - Heuristic example: Naive Manhattan Distance (sum of Manhattan distances from each box to nearest goal)
% - Baseline: Greedy matching + mobility penalty
% - Hyperparameters: 10 iterations, maximum 30,000 search steps
% - Training on 20 Sokoban levels
% - Interesting finding: comparison of training on 20 levels vs selecting 20 from 50 levels (50-select-20)
% TODO: Add results figures and tables

\subsection{TextWorld Benchmark}
% Experiment setup:
% 1. Introduced two games: Coin and MapReader
% 2. Used train environments to collect gameplay traces
% 3. Selected the 3 hardest/failure cases for LLM heuristic optimization
% 4. Evaluated on larger, unseen test environments

% Behavior example:
% - Compared baseline heuristic vs. LLM-optimized heuristics for 4 phases
% - Base heuristic gets stuck in local loops
% - Optimized heuristic encourages global exploration

% Optimization results (Simulation heuristic difference):
% - Baseline: uniform random rollout
% - Optimized: score-guided rollout that rewards task actions and goal progress
% - Simulation optimization improves test correctness → generalized heuristic

\begin{table*}[t]
  \caption{Results across configurations on train/test splits. Correct is \#correct/\#total; $\Delta$ is the difference in Correct vs. the baseline on the same split; Ret and Steps are the mean return and mean steps.}
  \label{tab:textworld-results}
  \centering
  \small
  \setlength{\tabcolsep}{4pt}
  \resizebox{\textwidth}{!}{%
  \begin{tabular}{lcccc|cccc}
    \toprule
    Configuration & Train Correct & Train $\Delta$ & Train Ret & Train Steps & Test Correct & Test $\Delta$ & Test Ret & Test Steps \\
    \midrule
    baseline & 62/72 & 0 & 0.861 & 9.9 & 55/90 & 0 & 0.606 & 14.4 \\
    \midrule
    + selection & 63/72 & 1 & 0.875 & 12.8 & 75/90 & 20 & 0.833 & 15.8 \\
    + expansion & 62/72 & 0 & 0.819 & 10.9 & 55/90 & 0 & 0.478 & 11.3 \\
    + simulation & \textbf{\underline{72/72}} & \textbf{\underline{10}} & \textbf{\underline{1.000}} & 41.0 & \textbf{\underline{90/90}} & \textbf{\underline{35}} & \textbf{\underline{1.000}} & 45.1 \\
    + backpropagation & 64/72 & 2 & 0.889 & \textbf{\underline{8.5}} & 60/90 & 5 & 0.656 & \textbf{\underline{10.2}} \\
    + full & 70/72 & 8 & 0.972 & 14.9 & 89/90 & 34 & 0.989 & 16.2 \\
    \bottomrule
  \end{tabular}%
  }
\end{table*}

\subsection{Heuristic Generalizability}
% - LLM-MCTS pass rate: 60~90 (full) out of 90
% - DQN pass rate: 27~33 (optimized) out of 90 (test blind — not trained on test environment)
% - DQN techniques used: Boltzmann exploration, success replay buffer, stagnation threshold, early stopping, step cost penalty
% - Key finding: LLM-generated heuristics generalize better than DQN to unseen test environments

\subsection{Zork}
% Environment complexity (vs. TextWorld):
% - Up to 22 interconnected rooms (vs. simple grid), with surface + underground layers
% - 6 treasures to collect AND deposit back at base (vs. pick up 1 coin)
% - Multi-step dependencies: get lantern → light it → go underground → kill troll → access treasures → return to deposit
% - Instant death states: entering dark rooms without lantern
% - Resource management: limited lantern fuel
% - State space orders of magnitude larger (room × inventory × items × fuel × troll × lantern × trophy case)
% TODO: Add Zork results figure


\section{Limitations}
\label{sec:limitations}

% - Were not able to meaningfully evaluate long-horizon behavior of tool quality
% - Open question: After ~1000 code iterations, does the tool quality collapse, plateau, or see eureka moments of continuous improvement?
% - Iterative task proposal can get stuck in parameter tuning without proposing new ideas; task proposal could use more research
% - MCTS heuristic tools show some transferability on same game with different setups
% - Lacks investigation on cross-game transferable tools
% - Evaluation limited to simple environments
% - Agent demonstrated improvements on multiple simple games, but lacked a larger, more impressive game environment to evaluate high-impact improvements


\section{Future Work}
\label{sec:future}

% Broader Agent Scope:
% - We limited LLM agent's scope of read/write access to only specific tool files within specific RL frameworks (MCTS) for feasible evaluation
% - The LLM agent might perform better with elevated scope — access to full tool folders, freedom to switch RL frameworks that interact with the environment, or freedom to modify its own prompt engine/LLM workflow

% Task Proposal Effectiveness:
% - Further research needed on improving the quality and diversity of task proposals

% Generalizing Beyond White-Box Games:
% - The game environments we evaluated are whitebox with explicit code endpoints for interaction
% - Potentially allow LLM agent to figure out / implement interaction frameworks with blackbox games
% - For example, using screenshots to learn game interaction but implementing tools to later build encapsulated movesets with mouse clicks


\section{Conclusion}
\label{sec:conclusion}

% TODO: Write conclusion summarizing:
% - We presented a self-evolving game agent that uses LLM to generate and improve heuristics for MCTS and A*
% - The system collects game traces, analyzes feedback, proposes tasks, and generates pluggable heuristics iteratively
% - Evaluated on Sokoban (A*), TextWorld (MCTS), and Zork environments
% - LLM-generated heuristics showed improvement over baselines and demonstrated generalizability to unseen environments
% - The approach offers a cost-effective alternative to training deep neural networks or using LLM as an online agent


\section*{References}

{\small

[1] Browne, C.~B., Powley, E., Whitehouse, D., Lucas, S.~M., Cowling, P.~I., Rohlfshagen, P., Tavener, S., Perez, D., Samothrakis, S., \& Colton, S. (2012). A Survey of Monte Carlo Tree Search Methods. \textit{IEEE Transactions on Computational Intelligence and AI in Games}, 4(1), 1--43.

[2] Shi, Z., Fang, M., \& Chen, L. (2025). Monte Carlo Planning with Large Language Model for Text-Based Game Agents. \textit{arXiv preprint arXiv:2504.16855}.

[3] Li, H., et al. (2026). Yunjue Agent Tech Report: A Fully Reproducible, Zero-Start In-Situ Self-Evolving Agent System for Open-Ended Tasks. \textit{arXiv preprint arXiv:2601.18226}.

}


%%%%%%%%%%%%%%%%%%%%%%%%%%%%%%%%%%%%%%%%%%%%%%%%%%%%%%%%%%%%

\appendix

\section{Technical Appendix}
% TODO: Add any supplementary material, additional results, or proofs here.

%%%%%%%%%%%%%%%%%%%%%%%%%%%%%%%%%%%%%%%%%%%%%%%%%%%%%%%%%%%%


\end{document}